% !TEX program = lualatex
\documentclass[aspectratio=43,professionalfonts,handout]{beamer}
\usepackage{beamer_template}

\newcommand{\LectureNum}{2}
\newcommand{\LectureTitle}{相似性と移動法則}
\newcommand{\TextPages}{p.\,50-53}
\newcommand{\CourseName}{応用数学}
\renewcommand{\TermName}{前期}

\begin{document}

\begin{frame}{\CourseName\quad 第\LectureNum 回「\LectureTitle」}
  {\small \TermName\quad \TeacherName\quad ／\quad 教科書 \TextPages}

  \textbf{目標}：ラプラス変換の相似性・移動法則を理解し、より複雑な関数の変換ができる

  \begin{block}{概要}
  \begin{itemize}
    \item 相似性
    \item 像関数の移動法則
    \item 原関数の移動法則
  \end{itemize}
  \end{block}
\end{frame}

\begin{frame}{相似性}
  \begin{block}{}
  \begin{itemize}
    \item $\mathcal{L}[f(at)] = (1/a)F(s/a)$
    \item 証明：置換積分 $at=\tau$
    \item 応用：$\mathcal{L}[\sin \omega t] = \omega/(s^{2}+\omega^{2})$
  \end{itemize}
  \end{block}
\end{frame}

\begin{frame}{像関数の移動法則}
  \begin{block}{}
  \begin{itemize}
    \item $\mathcal{L}[e^{\alpha t}f(t)] = F(s-\alpha)$
    \item s を $s-\alpha$ にシフト
    \item 例：$\mathcal{L}[e^{\alpha t}\sin \beta t] = \beta/((s-\alpha)^{2}+\beta^{2})$
    \item 例：$\mathcal{L}[e^{\alpha t}\cos \beta t] = (s-\alpha)/((s-\alpha)^{2}+\beta^{2})$
  \end{itemize}
  \end{block}
\end{frame}

\begin{frame}{原関数の移動法則}
  \begin{block}{}
  \begin{itemize}
    \item $\mathcal{L}[f(t-\mu)U(t-\mu)] = e^{-\mu s}F(s)$
    \item 時間軸のシフト
    \item グラフのイメージ
  \end{itemize}
  \end{block}
\end{frame}

\begin{frame}{演習}
  \begin{exampleblock}{自分で解いてみよう}
  \begin{itemize}
    \item 問7〜12より選題
  \end{itemize}
  \end{exampleblock}
\end{frame}

\begin{frame}{まとめ}
  \begin{itemize}
    \item 3つの法則の使い分け
    \item 次回：微分法則と積分法則
  \end{itemize}
\end{frame}

\end{document}
